%==============================================================================
\section{Introduction} \label{sec:introducao}
%==============================================================================

Infrared thermography captures radiation emitted by a surface, which depends on its temperature, and therefore measures skin temperature without contact or external illumination \cite{Nowakowski2025}.
A temperature measurement, however, is only as valid as the region it is averaged over.
Facial temperature is not uniformly distributed across regions that differ in vascular supply, tissue thickness, and sensitivity to ambient conditions: the inner corner of the eye correlates strongly with core temperature, whereas the forehead and cheeks are more susceptible to environmental artifacts \cite{Ioannou2014,Gioia2024}.
Isolating each anatomical region is thus a methodological prerequisite for downstream measurement, and because manual delineation is impracticable frame by frame, it must be automated by semantic segmentation.

Multi-region segmentation of thermal faces remains underexplored: most work targets face detection, landmark localization, or a single region of interest \cite{Vega2025,Chu2019,Jerbi2024}, and thermal segmentation as a whole is a fragmented subfield with heterogeneous evaluation protocols \cite{Kutuk2022}.
More importantly, architecture comparisons in this setting are fragile because of four recurring practices.
(i)~\textbf{Identity leakage}: when frames of the same subject appear in training and evaluation, metrics reward memorization of faces rather than generalization.
(ii)~\textbf{Identical hyperparameters} across convolutional and attention-based families implicitly adopt one family's recipe as the standard and handicap the other.
(iii)~\textbf{Undeclared training budgets} let an architecture look superior merely because it was trained longer.
(iv)~\textbf{Cost measured under unequal conditions}, such as different batch sizes or unsynchronized GPU timers, makes efficiency rankings incomparable.
Isensee et al.~\cite{Isensee2024} argue that much of the progress reported in medical segmentation stems from such validation practices rather than from architectural innovation.

This paper presents \textbf{UBench}, a reproducible benchmarking framework that addresses these four methodological issues by design and uses it to compare U-Net \cite{Ronneberger2015}, TransUNet \cite{Chen2021}, and \swinp{}, a simplified Swin-based baseline, on 10 facial classes of the Charlotte-ThermalFace dataset \cite{Ashrafi2022}.
We answer three research questions:
\textbf{RQ1}: Which architecture achieves the highest accuracy, and is the difference statistically distinguishable from experimental variability?
\textbf{RQ2}: How does per-class performance relate to region size and to the fidelity of the derived reference masks?
\textbf{RQ3}: What is the relationship between accuracy and computational cost under equalized conditions, when cost is measured?

The contributions are:
\begin{itemize}
\item an open, end-to-end benchmarking framework, from radiometric preprocessing to paired statistical analysis, in which all architectures share one training loop, one metric implementation, and one cost-measurement procedure;
\item a benchmark of three architectures under subject-disjoint cross-validation, family-specific recipes, an identical declared budget, and subject-paired significance tests;
\item the empirical finding that, in this domain, the variability across validation partitions exceeds the difference between the best and worst architectures, together with a sensitivity analysis showing that only a \emph{partial} ordering survives the exclusion of classes affected by label defects;
\item a measurement-methodology account, drawn from a 32-item audit of the pipeline, of how silent instrumentation defects reversed the memory ranking twice and corrupted labels without raising any errors.
\end{itemize}

Section~\ref{sec:trab_relac} reviews related work; Section~\ref{sec:metodologia} describes the framework and protocol; Section~\ref{sec:resultados} reports the results; Section~\ref{sec:discussao} discusses them; and Section~\ref{sec:conclusao} concludes. 